\chapter{Capítulo Sexto - Presupuesto y evaluación de costes}


En este capítulo se realiza una evaluación económica del proyecto TankGo desde dos perspectivas
complementarias. Por un lado, se estima el \textbf{coste teórico} del proyecto, es decir, lo que hubiera
supuesto contratar en el mercado el trabajo técnico realizado. Por otro lado, se detalla el
\textbf{coste real} incurrido durante el desarrollo, que en gran medida se ha mantenido en valores
mínimos gracias al uso estratégico de herramientas gratuitas y niveles de uso libres de coste
(\textit{free tiers}) de los principales proveedores cloud.

La diferencia entre ambas magnitudes es en sí misma un resultado relevante del proyecto: demuestra
que es posible construir y desplegar una plataforma de microservicios cloud-native de complejidad
real con una inversión económica muy reducida durante la fase de desarrollo y prototipado.

% ------------------------------------------------------------
\section{Costes laborales}
% ------------------------------------------------------------

El coste laboral constituye la partida más significativa del proyecto, dado que TankGo ha sido
desarrollado en su totalidad por un único estudiante. A lo largo del proyecto, el autor ha desempeñado
simultáneamente varios roles técnicos, cuyas horas se desglosan a continuación.

\subsection{Roles y tarifa de referencia}

Para estimar el coste de mercado del trabajo realizado se emplea una tarifa de referencia de
\textbf{30 €/hora}, correspondiente al perfil de un Ingeniero de Software Junior con conocimientos en
desarrollo fullstack, arquitecturas cloud y aprendizaje automático. Esta cifra es coherente con las
tablas salariales publicadas para perfiles similares en el mercado español en 2025--2026
\cite{infojobs2025salarios}.

\subsection{Distribución de horas por rol}

La Tabla~\ref{tab:costes_laborales} recoge la distribución estimada de horas invertidas por rol técnico
durante los aproximadamente nueve meses de desarrollo del proyecto (octubre de 2025 a junio de 2026).
La estimación parte de los registros de actividad y los hitos del diagrama de Gantt presentado en el
Capítulo~4.

\begin{table}[H]
\centering
\caption{Distribución de horas por rol y coste laboral estimado}
\label{tab:costes_laborales}
\begin{tabular}{lrrc}
\hline
\textbf{Rol} & \textbf{Horas} & \textbf{Tarifa (€/h)} & \textbf{Subtotal (€)} \\
\hline
Arquitecto de software / diseño del sistema   & 80  & 30 & 2.400 \\
Desarrollador backend (Python/FastAPI)        & 120 & 30 & 3.600 \\
Desarrollador backend (Node.js/Fastify/Hono)  & 80  & 30 & 2.400 \\
Desarrollador frontend (React/TypeScript)     & 70  & 30 & 2.100 \\
Ingeniería de datos / ML (LightGBM, pipeline) & 60  & 30 & 1.800 \\
DevOps / cloud / CI-CD (GCP, Docker)          & 50  & 30 & 1.500 \\
Análisis de requisitos y documentación        & 40  & 30 & 1.200 \\
\hline
\textbf{Total estudiante}                     & \textbf{500} & 30 & \textbf{15.000} \\
\hline
\end{tabular}
\end{table}

\subsection{Coste del director del proyecto}

La normativa del PFG establece que debe incluirse también el coste del equipo de soporte,
concretamente el del director. Se estima una dedicación media de 1 hora semanal de reunión y
seguimiento a lo largo de 36 semanas de proyecto, más un período de revisión de la memoria. La
tarifa de referencia para un perfil de Profesor Titular es de 60~€/hora.

\begin{table}[H]
\centering
\caption{Coste del director del proyecto}
\label{tab:coste_director}
\begin{tabular}{lrrc}
\hline
\textbf{Concepto} & \textbf{Horas} & \textbf{Tarifa (€/h)} & \textbf{Subtotal (€)} \\
\hline
Reuniones y seguimiento semanal & 36 & 60 & 2.160 \\
Revisión de la memoria          & 10 & 60 &   600 \\
\hline
\textbf{Total director}         & \textbf{46} & 60 & \textbf{2.760} \\
\hline
\end{tabular}
\end{table}

% ------------------------------------------------------------
\section{Material, herramientas e infraestructura}
% ------------------------------------------------------------

\subsection{Hardware}

El único elemento hardware empleado ha sido el ordenador personal del autor, un equipo portátil
de uso general valorado en 1.200~€. Se aplica una amortización lineal a 4 años, con un uso
dedicado al proyecto de aproximadamente el 60\% del tiempo durante 9 meses:

%\[
%\text{Amortización} = \frac{1200}{48}\times 9 \times 0{,}60 \approx 135\,\text{€}
%\]

\subsection{Software y herramientas de desarrollo}

La práctica totalidad del software empleado en el proyecto es de código abierto o gratuito, lo que
ha supuesto una ventaja económica significativa. La Tabla~\ref{tab:sw_herramientas} recoge las
principales herramientas utilizadas.

\begin{table}[H]
\centering
\caption{Herramientas de software y su coste}
\label{tab:sw_herramientas}
\begin{tabular}{llr}
\hline
\textbf{Herramienta} & \textbf{Uso} & \textbf{Coste} \\
\hline
Python, Node.js, React      & Lenguajes y frameworks principales   & Gratuito \\
Docker Desktop              & Contenedores y orquestación local    & Gratuito \\
Visual Studio Code          & Editor de código                     & Gratuito \\
Git / GitHub                & Control de versiones y CI/CD         & Gratuito \\
Overleaf                    & Edición de la memoria en \LaTeX{}    & Gratuito \\
\hline
\textbf{Total software}     &                                      & \textbf{0 €} \\
\hline
\end{tabular}
\end{table}

\subsection{Infraestructura cloud — coste real durante el desarrollo}
\label{subsec:infra_real}

Uno de los resultados más destacados desde el punto de vista económico es que la infraestructura
cloud utilizada durante todo el ciclo de desarrollo se ha mantenido dentro de los niveles gratuitos
ofrecidos por los proveedores. La Tabla~\ref{tab:infra_real} detalla los servicios empleados y su
coste efectivo.

\begin{table}[H]
\centering
\caption{Infraestructura cloud — coste real durante el desarrollo}
\label{tab:infra_real}
\begin{tabular}{lp{5.5cm}r}
\hline
\textbf{Servicio} & \textbf{Descripción del uso} & \textbf{Coste real} \\
\hline
Neon (PostgreSQL serverless) &
  Base de datos compartida para gasolineras, usuarios e historial de precios.
  El plan gratuito incluye 0,5 GiB de almacenamiento y cómputo bajo demanda. &
  0 € \\
Google Cloud Run &
  Despliegue de los microservicios (gateway, gasolineras, usuarios, recomendación).
  El tráfico del proyecto académico no ha superado los umbrales de facturación
  en ningún mes. &
  0 € \\
Google Cloud Storage &
  Almacenamiento de snapshots Parquet para el pipeline de ML y artefactos
  de modelos LightGBM. Uso inferior a 5 GiB, cubierto por el nivel gratuito. &
  0 € \\
Google Cloud Build &
  Pipelines de CI/CD (7 pipelines independientes). Los primeros 120 minutos
  de construcción al día son gratuitos, suficientes para el ritmo de despliegues
  del proyecto. &
  0 € \\
Artifact Registry &
  Registro de imágenes Docker. El nivel gratuito incluye 0,5 GiB por región,
  suficiente para imágenes de desarrollo. &
  $\approx$ 0 € \\
Google AI Studio (Gemini 2.5 Flash) &
  Asistente de voz con pipeline STT-LLM-TTS. Utilizado con la capa gratuita
  de Google AI Studio durante el desarrollo y pruebas. &
  0 € \\
Dominio \texttt{tankgo.dev} &
  Registro anual del dominio. &
  $\approx$ 11 € \\
\hline
\textbf{Total infraestructura} & & \textbf{$\approx$ 11 €} \\
\hline
\end{tabular}
\end{table}

Es importante señalar que el coste efectivamente nulo de los servicios cloud se debe al carácter
académico del proyecto y a la escala de tráfico durante la fase de desarrollo. En un escenario de
producción con usuarios reales, estos costes se incrementarían de forma proporcional al uso, tal
como se analiza en la sección~\ref{subsec:infra_prod}.

% ------------------------------------------------------------
\section{Coste total estimado del proyecto}
% ------------------------------------------------------------

La Tabla~\ref{tab:coste_total} consolida todas las partidas descritas en las secciones anteriores
para obtener el coste total estimado del proyecto si este se hubiera contratado en el mercado
bajo condiciones laborales estándar.

\begin{table}[H]
\centering
\caption{Coste total estimado del proyecto}
\label{tab:coste_total}
\begin{tabular}{lr}
\hline
\textbf{Partida} & \textbf{Importe (€)} \\
\hline
Costes laborales — estudiante (500 h $\times$ 30 €/h) & 15.000 \\
Costes laborales — director (46 h $\times$ 60 €/h)    &  2.760 \\
Amortización de hardware                               &    135 \\
Software y herramientas                                &      0 \\
Infraestructura cloud (dominio incluido)               &     11 \\
\hline
\textbf{Total estimado}                                & \textbf{17.906 €} \\
\hline
\end{tabular}
\end{table}

% ------------------------------------------------------------
\section{Estimación del coste en producción}
\label{subsec:infra_prod}
% ------------------------------------------------------------

Con el fin de contextualizar la viabilidad económica del proyecto más allá de la fase académica,
se ofrece a continuación una estimación orientativa del coste mensual de infraestructura que
tendría TankGo en un escenario de producción real con carga moderada.

\subsection{Cloud Run — microservicios}

Google Cloud Run factura según el consumo de CPU, memoria y peticiones. Para un servicio
desplegado en \texttt{europe-west1} con una carga de 500.000 peticiones mensuales y una
latencia media de 300~ms por petición, el coste estimado se sitúa en torno a \textbf{5--15~€/mes}
por servicio, según la calculadora oficial de Google Cloud \cite{gcp_run_pricing}. Con los siete
servicios del proyecto (de los cuales algunos tienen tráfico muy bajo), el coste total de Cloud Run
se estima en \textbf{20--50~€/mes}.

\subsection{Neon PostgreSQL}

Superado el plan gratuito, el plan Pro de Neon parte de \textbf{19~USD/mes} e incluye hasta
10~GiB de almacenamiento y cómputo escalable. Para un proyecto con el volumen de datos de
TankGo (gasolineras, historial de precios, usuarios), este plan sería suficiente en una primera
fase de producción.

\subsection{Google Gemini 2.5 Flash — asistente de voz}

El modelo Gemini 2.5 Flash, empleado en el pipeline de voz (STT, razonamiento con herramientas
y TTS), tiene un precio de \textbf{0,30~USD por millón de tokens de entrada} y \textbf{2,50~USD por
millón de tokens de salida} en texto \cite{gemini_pricing_2025}. El audio de entrada se factura
a \textbf{1,00~USD por millón de tokens de audio}. Dado que cada interacción de voz consume
un volumen de tokens muy reducido (del orden de miles), el coste unitario por consulta es
prácticamente despreciable, estimándose en menos de \textbf{0,001~USD por interacción}.

\subsection{Estimación mensual consolidada en producción}

\begin{table}[H]
\centering
\caption{Estimación orientativa de costes mensuales en producción}
\label{tab:prod_cost}
\begin{tabular}{lr}
\hline
\textbf{Servicio} & \textbf{Coste mensual estimado} \\
\hline
Google Cloud Run (7 servicios)     & 20 -- 50 € \\
Neon PostgreSQL (plan Pro)         & $\approx$ 18 € \\
Google Cloud Storage               & $<$ 1 € \\
Gemini 2.5 Flash (uso moderado)    & 1 -- 5 € \\
Dominio y DNS                      & $\approx$ 1 € \\
\hline
\textbf{Total mensual estimado}    & \textbf{40 -- 75 €/mes} \\
\hline
\end{tabular}
\end{table}

Esta estimación pone de manifiesto que TankGo podría operar como servicio real con un coste
de infraestructura inferior a \textbf{1.000~€/año}, lo que refleja la eficiencia económica de las
arquitecturas serverless y de los modelos pay-as-you-go para proyectos de escala media.

% ------------------------------------------------------------
\section{Análisis y valoración económica}
% ------------------------------------------------------------

Del análisis anterior se extraen las siguientes conclusiones económicas relevantes:

En primer lugar, el coste estimado total del proyecto asciende a \textbf{17.906~€}, cifra que refleja
fielmente el volumen y la complejidad del trabajo técnico realizado: diseño de una arquitectura de
microservicios con siete componentes, desarrollo fullstack, integración de modelos de aprendizaje
automático, despliegue en Google Cloud y extensión con capacidades conversacionales basadas en
IA generativa.

En segundo lugar, el coste real incurrido ha sido prácticamente nulo en lo que respecta a
infraestructura, gracias a la selección cuidadosa de herramientas con niveles gratuitos generosos
(Neon, GitHub, Google AI Studio, Cloud Run dentro de los umbrales de prueba). El único gasto
efectivo ha sido el registro del dominio \texttt{tankgo.dev} por un importe de aproximadamente
11~€. Esta decisión de diseño, lejos de ser una casualidad, es resultado de priorizar tecnologías
cloud-native modernas cuyo modelo de negocio está específicamente orientado a facilitar el
desarrollo y la experimentación sin coste inicial.

Por último, la estimación de costes en producción —entre 40 y 75~€/mes— demuestra que el
proyecto es económicamente viable como servicio real, con una estructura de costes escalable y
predecible, característica deseable en cualquier producto software orientado a su eventual
comercialización o despliegue a mayor escala.